\documentclass[10pt,a4paper]{article}
\usepackage[utf8]{inputenc}
\usepackage[T1]{fontenc}
\usepackage{lmodern}
\usepackage[spanish]{babel}
\usepackage[margin=14mm,top=13mm,bottom=16mm]{geometry}
\usepackage{xcolor,enumitem,titlesec,tabularx,fancyhdr,lastpage}
\usepackage[hidelinks,unicode]{hyperref}
\definecolor{navy}{HTML}{1D405C}
\definecolor{muted}{HTML}{515B64}
\definecolor{rulecolor}{HTML}{D6DEE4}
\renewcommand{\familydefault}{\sfdefault}
\setlength{\parindent}{0pt}
\setlength{\parskip}{3pt}
\setlist[itemize]{leftmargin=12pt,itemsep=2pt,topsep=3pt,parsep=0pt}
\titleformat{\section}{\color{navy}\small\bfseries}{}{0em}{}[\color{rulecolor}\titlerule]
\titlespacing*{\section}{0pt}{11pt}{5pt}
\newcommand{\entry}[3]{\noindent\begin{tabularx}{\linewidth}{@{}Xr@{}}\textbf{#1}&{\small\color{muted}#2}\end{tabularx}\par{\small\color{muted}#3}\par}
\pagestyle{fancy}\fancyhf{}\renewcommand{\headrulewidth}{0pt}
\fancyfoot[L]{\scriptsize\color{muted}Alejandro Villarrubia · Technical Sales Engineer}
\fancyfoot[R]{\scriptsize\color{muted}\thepage\,/\,\pageref{LastPage}}
\hypersetup{pdftitle={Alejandro Villarrubia -- CV},pdfauthor={Alejandro Villarrubia}}
\begin{document}
{\fontsize{23}{26}\selectfont\bfseries\color{navy} ALEJANDRO VILLARRUBIA}\par
{\large Technical Sales Engineer · Software Engineer}\par
{\small\color{muted}+34 633 58 60 60 · \href{mailto:avillarrubia2@gmail.com}{avillarrubia2@gmail.com} · Palencia, España}\par
{\small\href{https://github.com/villarrubi}{github.com/villarrubi} · \href{https://www.linkedin.com/in/alejandro-villarrubia-garc\%C3\%ADa-3078892b9/}{LinkedIn}}\par
\section*{PERFIL PROFESIONAL}
Ingeniero informático con experiencia en desarrollo de software, integración hardware/software y coordinación entre ingeniería, ventas y cliente. Actualmente Technical Sales Engineer en ALTEN Spain para MAHLE, especializado en sistemas de propulsión para e-bike. Experiencia previa en backend con Java, MongoDB y SQL. Interés por construir sistemas desde cero, inteligencia artificial, algoritmos y proyectos técnicos propios.
\section*{EXPERIENCIA PROFESIONAL}
\entry{Technical Sales Engineer}{jul. 2026 -- presente}{ALTEN Spain · cliente: MAHLE | Palencia}
\begin{itemize}
\item Coordino la planificación y seguimiento de flotas de sistemas de propulsión para e-bike junto a Program y Product Management, gestionando recursos, entregables y documentación técnica.
\item Valido la compatibilidad técnica de cuadros de cliente (Customer Frame Analysis) con los sistemas MAHLE antes de producción.
\item Participo en test rides y configuración de motor maps (potencia y par), colaborando con Battery and Control Systems para evaluar motores y baterías.
\item Actúo como enlace entre ventas, ingeniería y cliente, preparando informes técnicos y gestionando muestras y documentación asociada.
\end{itemize}
\entry{Software Engineer}{oct. 2025 -- jul. 2026}{NTT DATA · España (equipo distribuido nacional)}
\begin{itemize}
\item Desarrollé de forma autónoma una nueva funcionalidad de producto desde cero a partir de una solución existente, participando en el ciclo completo de entrega en Agile/Scrum.
\item Refactoricé código, corregí errores y resolví incidencias de producción para mejorar mantenibilidad y reducir deuda técnica en un equipo de más de 20 ingenieros.
\item Participé en diseño, implementación, pruebas, depuración, revisión de código y despliegue.
\end{itemize}
{\small\color{muted} Java · SQL · MongoDB · Git · Agile/Scrum}\par\vspace{4pt}
\entry{Prácticas -- Desarrollo de software}{abr. 2025 -- jul. 2025}{REDALSA, S.A.}
\begin{itemize}
\item Diseñé e implementé desde cero un sistema de detección de objetos en vía mediante LiDAR, con medición de distancias, mapeo y cálculo en tiempo real de posibles colisiones y trayectorias.
\item Integré hardware y software de forma autónoma, desde captura y procesamiento de señal hasta visualización e interpretación de datos.
\end{itemize}
{\small\color{muted} Python · C · LiDAR · Sistemas embebidos · Adquisición de datos}\par\vspace{4pt}
\entry{Miembro activo de comisiones}{abr. 2024 -- oct. 2025}{RITSI · representación nacional de estudiantes de Informática}
\begin{itemize}
\item Participé en cuatro comisiones estatales con más de 30 universidades, incluyendo organización de congresos nacionales y mantenimiento de herramientas tecnológicas.
\end{itemize}
\section*{FORMACIÓN}
\entry{Grado en Ingeniería Informática}{2021 -- 2026}{Universidad Europea Miguel de Cervantes (UEMC)}\vspace{3pt}
\entry{Máster en Inteligencia Artificial e Innovación}{2025}{Microsoft · Founderz | Formación de posgrado}\vspace{3pt}
\entry{Bachillerato de Ciencias}{2015 -- 2021}{IES Jorge Manrique}\vspace{3pt}
\newpage
\section*{PROYECTOS SELECCIONADOS}
\entry{ChessBot -- Motor de ajedrez}{2026 -- presente}{Proyecto personal · C++20 · UCI · IA local}
\begin{itemize}
\item Motor propio capaz de jugar y analizar posiciones, con sistemas de búsqueda y evaluación y comunicación UCI con interfaces y otros motores.
\item Infraestructura de entrenamiento con autojuego, partidas contra otros motores y experimentación con evaluación posicional y redes NNUE.
\item Aplicación Windows e integración opcional de un LLM local para explicar cálculos del motor. Proyecto en entrenamiento y evaluación, sin Elo certificado.
\end{itemize}
\entry{Pip$\alpha$ -- Asistente local y controlador hardware}{2026 -- presente}{Proyecto personal · Python · ESP32-S3 · Windows · Swift}
\begin{itemize}
\item Agente local para automatizar aplicaciones, audio, estado del sistema y acciones de control desde un dispositivo externo, con comandos confirmables.
\item Integración hardware/software: pantalla táctil, voz local, comunicación USB cifrada y Wake-on-LAN. Cliente iOS experimental.
\end{itemize}
\entry{Detección de phishing en correo electrónico}{2026}{Trabajo de Fin de Grado · Python · scikit-learn · Streamlit}
\begin{itemize}
\item Análisis heurístico de 28 señales (cabeceras, URLs, HTML y SPF/DKIM/DMARC) y clasificador neuronal MLP entrenado con datasets reales.
\item Pipeline TF-IDF + MLP, interfaces web, extensión de Gmail, pruebas y persistencia del modelo; análisis en español e inglés.
\end{itemize}
\entry{AQUILAE -- Simulador Risk Imperio Romano}{2025}{Colaboración en TFM · Java · Swing · Maven}
\begin{itemize}
\item Diseñé e implementé el simulador completo, con mapa histórico vectorial de 28 provincias e interfaz Java Swing, como apoyo a un TFM de la Universidad de Valladolid.
\end{itemize}
\section*{HABILIDADES TÉCNICAS}
\textbf{Lenguajes:} Java, Python, C, C++, C\#, JavaScript, HTML, PHP, SQL.\par
\textbf{Datos y herramientas:} SQL, MongoDB, Git, Agile/Scrum, Office 365.\par
\textbf{Especialidades:} Algoritmos de búsqueda, UCI, IA/ML, integración hardware/software, LiDAR, sistemas embebidos.\par
\section*{CERTIFICACIONES E IDIOMAS}
B2 First -- Cambridge University · Certificado de Competencias Digitales (tuCertiCyL) · Monitor Autonómico de Ajedrez · Árbitro Autonómico de Ajedrez.\par
Español nativo · Inglés B2 (Cambridge First) · Alemán básico.
\section*{REPRESENTACIÓN Y COMUNIDAD}
\entry{Delegado universitario · EPS, Informática y curso}{oct. 2021 -- oct. 2025}{UEMC · Representación institucional e interlocución con dirección académica y decanato.}\vspace{4pt}
\entry{Representante de la UEMC en RITSI}{oct. 2023 -- nov. 2025}{Delegado oficial ante la sectorial nacional de Informática; comisiones y congresos.}\vspace{4pt}
\entry{Secretario · Club Palentino de Ajedrez}{mar. 2026 -- presente}{Dirección y organización de actividades de un club con más de 40 años de historia.}\vspace{4pt}
\entry{Creador y coordinador}{jun. 2024 -- presente}{Club de Lectura UEMC · Escuela de Ajedrez UEMC.}\vspace{4pt}
\end{document}
